
\section{Gaz molekulalarining o`rtacha tezliklari}

\begin{enumerate}

\item Ideal gazning absolyut temperaturasini necha marta ko‘targanda, molekulalarning o‘rtacha kvadratik tezligi 2 marta ortadi?
\begin{enumerate}[label=\Alph*)]
\item $\sqrt{2}$
\item 4
\item 2
\item 16
\end{enumerate}

\item Kislorod molekulasi issiqlik harakatining o‘rtacha kvadratik tezligi vodorod molekulasinikidan necha marta kichik?
\begin{enumerate}[label=\Alph*)]
\item 2
\item 3
\item 4
\item 8
\end{enumerate}

\item Agar azot gazining absolyut temperaturasi 4 marta ortsa, molekulalarining o‘rtacha kvadratik tezligi qanday o‘zgaradi?
\begin{enumerate}[label=\Alph*)]
\item 4 marta kamayadi
\item 4 marta ortadi
\item 2 marta kamayadi
\item 2 marta ortadi
\end{enumerate}

\item Ideal gazning absolyut harorati 2 marta kamayganda molekulalarining o‘rtacha kvadratik tezligi necha marta kamayadi?
\begin{enumerate}[label=\Alph*)]
\item $\sqrt{2}$
\item 2
\item $2\sqrt{2}$
\item 4
\end{enumerate}

\item Gaz temperaturasi 21 \% ga ortsa, molekulalarning o‘rtacha kvadratik tezligi necha foizga ortadi?
\begin{enumerate}[label=\Alph*)]
\item 10
\item 50
\item 20
\item 25
\end{enumerate}

\item Temperatura qanday bo‘lganda (K) azot molekulalarining o‘rtacha kvadratik tezligi 830 m/s bo‘ladi? $M_{Azot} = 28 \, g/mol$.
\begin{enumerate}[label=\Alph*)]
\item 744
\item 774
\item 374
\item 425
\end{enumerate}

\item $27^\circ C$ haroratli vodorod molekulalarining o‘rtacha kvadratik tezligini (km/s) va shu tezlikdagi bitta molekulaning impulsini toping ($kg \cdot m/s$). Vodorodning molyar massasi $\mu = 2,0 \, g/mol$.
\begin{enumerate}[label=\Alph*)]
\item 2,46; $6,4 \cdot 10^{-27}$
\item 1,93; $6,4 \cdot 10^{-24}$
\item 1,93; $6,4 \cdot 10^{-21}$
\item 3,54; $6,4 \cdot 10^{-21}$
\end{enumerate}

\item Havoni bir jinsli gaz deb hisoblang, havoda muallaq turgan $1,74 \cdot 10^{-12} \, kg$ massali chang zarrasining o‘rtacha kvadratik tezligi havo molekulalari harakatining o‘rtacha kvadratik tezligidan necha marta kichik ekanini toping. $M_{havo} = 29 \, g/mol$.
\begin{enumerate}[label=\Alph*)]
\item $6 \cdot 10^4$ marta
\item $6 \cdot 10^6$ marta
\item $3 \cdot 10^6$ marta
\item $2 \cdot 10^6$ marta
\end{enumerate}

\item Ideal gazning mutlaq temperaturasi 3 marta ortganida, molekulalarining o‘rtacha kvadratik tezligi necha marta ortadi?
\begin{enumerate}[label=\Alph*)]
\item 9
\item $3\sqrt{3}$
\item 3
\item $\sqrt{3}$
\end{enumerate}

\item Atmosferadagi qaysi gaz molekulalari tezroq harakat qiladi?
\begin{enumerate}[label=\Alph*)]
\item $N_2$
\item $O_2$
\item $H_2$
\item $CO_2$
\end{enumerate}

\item Gazning temperaturasi $10^\circ C$ dan $20^\circ C$ gacha ortdi. Bunda molekulalar issiqlik harakatining o‘rtacha kvadratik tezligi qanday o‘zgaradi?
\begin{enumerate}[label=\Alph*)]
\item bir oz ortadi
\item 2 marta ortadi
\item o‘zgarmaydi
\item $\sqrt{2}$ marta ortadi
\end{enumerate}

\item Qanday temperaturada (K) vodorod molekulalarining o‘rtacha kvadratik tezligi 2 km/s bo‘ladi? Vodorodning molyar massasi $0,002 \, kg/mol$.
\begin{enumerate}[label=\Alph*)]
\item 273
\item 300
\item 350
\item 321
\end{enumerate}

\item Germetik berk idishdagi gaz molekulalarining kvadratik tezligi 2 marta ortishi uchun uning harorati $250^\circ$ dan qanchagacha ko‘tarilishi kerak?
\begin{enumerate}[label=\Alph*)]
\item $473^\circ$
\item $773^\circ$
\item $1819^\circ$
\item $1350^\circ$
\end{enumerate}

\item Gaz molekulalarining o‘rtacha kvadratik tezligi 600 m/s dan 750 m/s ga oshganda uning temperaturasi 30 K ortdi. O‘rtacha kvadratik tezligi 750 m/s dan 900 m/s ga ortsa, uning temperaturasi qancha ortadi (K)?
\begin{enumerate}[label=\Alph*)]
\item 27,3
\item 40,1
\item 30
\item 36,7
\end{enumerate}

\item Gaz molekulalarining o‘rtacha kvadratik tezligi 500 m/s dan 620 m/s ga oshganda uning temperaturasi 70 K ortdi. O‘rtacha kvadratik tezligi 620 m/s dan 800 m/s ga ortsa uning temperaturasi qancha ortadi (K)?
\begin{enumerate}[label=\Alph*)]
\item 77,3
\item 80,1
\item 93,2
\item 133,1
\end{enumerate}

\item 173 K temperaturadagi vodorod molekulalarining o‘rtacha kvadratik tezligini (m/s) va ilgarilanma harakat energiyasini hisoblang (J). $M_H = 2 \, g/mol$.
\begin{enumerate}[label=\Alph*)]
\item 2350, $1,8 \cdot 10^{-21}$
\item 1468, $3,6 \cdot 10^{-21}$
\item 1465, $3,2 \cdot 10^{-19}$
\item 2523, $3,2 \cdot 10^{-21}$
\end{enumerate}

\item Normal sharoitda bitta gaz molekulasi ilgarilanma harakatining o‘rtacha kinetik energiyasini (J) hisoblang.
\begin{enumerate}[label=\Alph*)]
\item $7,5 \cdot 10^{-20}$
\item $5,65 \cdot 10^{-21}$
\item $1,25 \cdot 10^{-21}$
\item $5,7 \cdot 10^{-20}$
\end{enumerate}

\item 273 K temperaturadagi vodorod molekulalarining o‘rtacha kvadratik tezligini (m/s) va ilgarilanma harakat energiyasini hisoblang (J). $M_H = 2 \, g/mol$.
\begin{enumerate}[label=\Alph*)]
\item 1840, $4,5 \cdot 10^{-21}$
\item 1844, $5,6 \cdot 10^{-21}$
\item 1205, $4,5 \cdot 10^{-21}$
\item 7520, $5,6 \cdot 10^{-21}$
\end{enumerate}

\item Shtern tajribasida tashqi silindrning ichki sirtida paydo bo‘ladigan kumush poloska (yo‘l) chaplashgan bo‘lib chiqadi. Bundan qanday xulosa chiqarish mumkin?
\begin{enumerate}[label=\Alph*)]
\item Molekulalarning harakat tezliklari turlicha
\item Molekulalarning harakat tezliklari bir xil
\item Molekulalarning harakat tezliklariga bog‘liq emas
\item Barcha javoblar noto'g'ri
\end{enumerate}

\item Agar Shtern tajribasida asbobning aylanish chastotasi $150 \, s^{-1}$ bo‘lgan kumush molekulasining burchak siljishi $5,4^\circ$ ni tashkil etsa, bunda kumush bug‘ining molekulasi qanday tezlikka ega bo‘lgan (m/s)? Ichki va tashqi silindrlar orasidagi masofa 2 sm ga teng.
\begin{enumerate}[label=\Alph*)]
\item 500
\item 100
\item 400
\item 200
\end{enumerate}

\end{enumerate}



\begin{table}[!ht]
\centering
\begin{tabular}{|c|c|c|c|c|c|c|c|c|c|}
\hline
№ & J. & № & J. & № & J. & № & J. & № & J. \\ \hline
1 & B & 5 & A & 9 & D & 13 & C & 17 & B \\ \hline
2 & C & 6 & B & 10 & C & 14 & D & 18 & B \\ \hline
3 & D & 7 & B & 11 & A & 15 & D & 19 & A \\ \hline
4 & A & 8 & B & 12 & D & 16 & B & 20 & D \\ \hline
\end{tabular}
\end{table}